\section{Problem Formulation}
\label{sec:problem}

We formalize LLM routing as an online decision-making problem under budget constraints, where the environment is non-stationary with respect to both model quality (reward drift) and model pricing (cost drift). 

\subsection{Contextual Bandit Routing}

The key premise of contextual routing is that \emph{no single model is universally best}---the optimal choice depends on the prompt.
LLMs exhibit heterogeneous capabilities across prompt types: a budget model may handle straightforward factual look-ups as well as a frontier API, while only the frontier model produces acceptable responses on complex multi-step reasoning or nuanced creative writing.
A context-free router that ignores prompt content must commit to a single, fixed traffic distribution (e.g., 70\% budget, 25\% mid-tier, 5\% frontier) regardless of prompt difficulty.
Under a budget constraint, this is wasteful: frontier calls are spent on easy prompts where a cheap model would suffice, while hard prompts that justify the $500{\times}$ cost premium receive frontier service no more often than easy ones.
Conditioning on the prompt allows the router to learn \emph{which} prompt types each model excels at, reserving expensive calls for the prompts where they yield the largest quality gain and routing the rest to cheaper alternatives.
This selective spending is what enables a contextual router to trace the quality--cost Pareto frontier---and it is impossible without a representation of the prompt that captures the task-level distinctions relevant to model selection.

Let $t \in \{1, 2, \ldots, T\}$ denote a discrete sequence of user requests (prompts). At each step $t$, the router receives a context vector $x_t \in \mathbb{R}^d$ representing the prompt, derived from a pre-trained embedding model (e.g., SentenceTransformers). The router must select a single language model $a_t$ from a discrete portfolio of available models $\mathcal{A} = \{1, \dots, K\}$. 

Upon executing the prompt on model $a_t$, the system observes a scalar reward $r_t \in [0, 1]$ indicating the quality of the response. This reward may be derived from an automated LLM-as-a-judge, explicit user feedback (e.g., thumbs up/down), or an implicit task completion metric. Crucially, the router operates under \emph{bandit feedback}: it only observes the reward for the chosen model $a_t$ and does not know what the counterfactual rewards would have been for the unselected models.

The primary objective is to learn a routing policy $\pi$ that maximizes the expected cumulative reward:
\begin{equation}
    \max_{\pi} \mathbb{E} \left[ \sum_{t=1}^T r_{t,\pi(x_t)} \right]
\end{equation}
In words: the router aims to make the best model choice on every request so that, summed over all $T$ requests, the total quality is as high as possible---even though it never observes what would have happened had it chosen differently.

\subsection{The Budget Constraint}

In production, each model $a \in \mathcal{A}$ incurs a known, deterministic inference cost $c_a(x_t)$ (typically priced per input and output token). Because the models in a portfolio often span orders of magnitude in price---from open-weights models costing fractions of a cent to frontier APIs costing several cents per request---the router must strictly manage aggregate spend. 

We formulate this as a \emph{Contextual Bandit with Knapsacks} (CBwK)~\cite{agrawal2014bandits} problem. The system is given a target average budget $B$ (e.g., ``average \$0.005 per req''). The policy $\pi$ must maximize cumulative reward subject to the constraint that the long-run average cost does not exceed $B$:
\begin{equation}
    \limsup_{T \to \infty} \frac{1}{T} \sum_{t=1}^T c_{\pi(x_t)}(x_t) \leq B
\end{equation}
In plain terms: averaged over many requests, the cost per request must stay at or below the target $B$.  The router may occasionally splurge on an expensive frontier model, but must offset that with cheaper choices elsewhere so the running average never exceeds the budget.

The standard approach to this constraint in the routing literature is to scalarize the reward using a fixed penalty weight $\lambda > 0$, optimizing a penalized utility $u_{a,t} = r_{a,t} - \lambda c_a(x_t)$~\cite{ong2024routellm}. However, selecting $\lambda$ requires extensive offline simulation on historical data to find the weight that happens to yield an average cost of $B$. If the distribution of prompts shifts, or if the portfolio changes, the realized cost will drift away from $B$, requiring manual recalibration.

\subsection{Non-Stationarity}

Real-world LLM deployments are fundamentally non-stationary. We define two specific types of environment drift that a production router must handle autonomously:

\paragraph{Reward Drift (Quality Shifts)}
LLM providers silently update their models behind APIs, or internal teams push new versions of fine-tuned open-weight models. Formally, the conditional expected reward function $f_a(x) = \mathbb{E}[r_t \mid x, a]$ changes at some unknown time $\tau$, such that $f_a^{(t < \tau)}(x) \neq f_a^{(t \geq \tau)}(x)$.  Put simply, the expected quality of a model's response to a given prompt is one value before the update and a different value after it---and the router does not know when the change occurred.  A router must detect this shift and rapidly unlearn its historical estimates for model $a$, avoiding a stale policy that continues routing traffic to a degraded model.

\paragraph{Cost Drift (Price Drops)}
API providers frequently reduce inference prices to remain competitive~\cite{thompson2026timeline}. Formally, the cost function $c_a(x)$ decreases at time $\tau$. If the router uses a fixed, offline-tuned penalty $\lambda$, it will continue to apply the same relative penalty to model $a$. While this mechanically reduces the total spend (pocketing the savings), it fails to recognize that model $a$ is now a relatively better value. An optimal router should dynamically exploit the newly freed budget to route more challenging queries to the newly-cheaper, higher-quality model, maximizing overall application quality while remaining precisely at the budget limit $B$.